% Figure: why a coordinate pair flips with probability theta/pi (supplement, Section S1.6). Drawn by hand; q at angle 0,
% q' at theta = 50 degrees; g flips the pair when it points into (90, 90+theta) or (270, 270+theta) degrees.
\begin{figure}[!htbp]
\centering
\begin{tikzpicture}[>=Stealth, scale=1.15]
\def\th{50}
\draw[black!30] (0,0) circle (2);
% the two arcs of directions of g that separate q and q'
\draw[red!70!black, line width=3pt] (90:2) arc (90:90+\th:2);
\draw[red!70!black, line width=3pt] (270:2) arc (270:270+\th:2);
% q and q'
\draw[->, very thick, blue!70!black] (0,0) -- (0:1.7) node[right] {$q$};
\draw[->, very thick, blue!70!black] (0,0) -- (\th:1.7) node[above right] {$q'$};
\draw (0:0.55) arc (0:\th:0.55);
\node at (\th/2:0.8) {$\theta$};
% one direction of g inside an arc, and the line orthogonal to it
\draw[->, thick, dashed] (0,0) -- (115:1.95) node[above left] {$g$};
\draw[dotted, very thick] (25:2.3) -- (205:2.3);
\node[font=\footnotesize, align=left, anchor=west, red!70!black] at (2.5,1.2) {$g$ in a red arc:\\the pair flips};
\node[font=\footnotesize, align=left, anchor=west] at (2.5,-0.3) {dotted: the line\\orthogonal to $g$};
\end{tikzpicture}
\caption{\textbf{Why a pair of coordinates flips with probability $\theta/\pi$.} The plane spanned by the two feature vectors $q$ and $q'$, which form the angle $\theta$. The projection of the random direction $g$ onto this plane points in a uniformly random direction. The pair is ordered differently for $q$ and $q'$ exactly when the line orthogonal to $g$ (dotted) separates $q$ from $q'$. That happens when $g$ points into one of the two red arcs. The arcs have total angle $2\theta$ out of $2\pi$, so the pair flips with probability $\theta/\pi$. The dashed arrow shows one such $g$.}
\label{fig:s-angle}
\end{figure}
